% Options for packages loaded elsewhere
\PassOptionsToPackage{unicode}{hyperref}
\PassOptionsToPackage{hyphens}{url}
\documentclass[
  american,
  11pt,
  letterpaper,
]{article}
\usepackage{xcolor}
\usepackage{amsmath,amssymb}
\setcounter{secnumdepth}{5}
\usepackage{iftex}
\ifPDFTeX
  \usepackage[T1]{fontenc}
  \usepackage[utf8]{inputenc}
  \usepackage{textcomp} % provide euro and other symbols
\else % if luatex or xetex
  \usepackage{unicode-math} % this also loads fontspec
  \defaultfontfeatures{Scale=MatchLowercase}
  \defaultfontfeatures[\rmfamily]{Ligatures=TeX,Scale=1}
\fi
\usepackage{lmodern}
\ifPDFTeX\else
  % xetex/luatex font selection
\fi
% Use upquote if available, for straight quotes in verbatim environments
\IfFileExists{upquote.sty}{\usepackage{upquote}}{}
\IfFileExists{microtype.sty}{% use microtype if available
  \usepackage[]{microtype}
  \UseMicrotypeSet[protrusion]{basicmath} % disable protrusion for tt fonts
}{}
\makeatletter
\@ifundefined{KOMAClassName}{% if non-KOMA class
  \IfFileExists{parskip.sty}{%
    \usepackage{parskip}
  }{% else
    \setlength{\parindent}{0pt}
    \setlength{\parskip}{6pt plus 2pt minus 1pt}}
}{% if KOMA class
  \KOMAoptions{parskip=half}}
\makeatother
\ifLuaTeX
\usepackage[bidi=basic,shorthands=off]{babel}
\else
\usepackage[bidi=default,shorthands=off]{babel}
\fi
\ifLuaTeX
  \usepackage{selnolig} % disable illegal ligatures
\fi
\setlength{\emergencystretch}{3em} % prevent overfull lines
\providecommand{\tightlist}{%
  \setlength{\itemsep}{0pt}\setlength{\parskip}{0pt}}
\usepackage[margin=1in]{geometry}
\usepackage{setspace}
\usepackage{fancyhdr}
\usepackage{titlesec}
\usepackage[hang,flushmargin,bottom]{footmisc}
\usepackage{booktabs}
\usepackage{longtable}
\usepackage{array}
\usepackage{enumitem}
\usepackage{caption}

\setstretch{1.08}
\setlength{\parindent}{0pt}
\setlength{\parskip}{0.7em}
\setlength{\footnotesep}{0.7em}
\interfootnotelinepenalty=10000
\emergencystretch=2em
\raggedbottom

\setlist[itemize]{leftmargin=1.5em, itemsep=0.45em, topsep=0.45em}
\setlist[enumerate]{leftmargin=1.8em, itemsep=0.45em, topsep=0.45em}

\pagestyle{fancy}
\fancyhf{}
\fancyhead[L]{\nouppercase{\leftmark}}
\fancyhead[R]{AI Empire Research Papers}
\fancyfoot[C]{\thepage}
\renewcommand{\headrulewidth}{0.3pt}

\setcounter{secnumdepth}{3}
\setcounter{tocdepth}{2}
\renewcommand{\contentsname}{Paper Roadmap}

\titleformat{\section}{\Large\bfseries}{\thesection.}{0.65em}{}
\titleformat{\subsection}{\large\bfseries}{\thesubsection.}{0.6em}{}
\titleformat{\subsubsection}{\normalsize\bfseries}{\thesubsubsection.}{0.55em}{}

\captionsetup{font=small, labelfont=bf}
\AtBeginEnvironment{longtable}{\small}

\AtBeginDocument{
  \hypersetup{
    colorlinks=true,
    linkcolor=black,
    citecolor=black,
    urlcolor=blue
  }
}
\usepackage{bookmark}
\IfFileExists{xurl.sty}{\usepackage{xurl}}{} % add URL line breaks if available
\urlstyle{same}
% fallback for those not using the hyperref driver hyperxmp:
\makeatletter
\@ifundefined{xmpquote}{\newcommand{\xmpquote}[1]{#1}}{}
\makeatother
\hypersetup{
  pdftitle={Chapter 06: AI as a Dual-Use Rhetorical Weapon},
  pdflang={en-US},
  hidelinks,
  pdfcreator={LaTeX via pandoc}}

\title{Chapter 06: AI as a Dual-Use Rhetorical Weapon}
\usepackage{etoolbox}
\makeatletter
\providecommand{\subtitle}[1]{% add subtitle to \maketitle
  \apptocmd{\@title}{\par {\large #1 \par}}{}{}
}
\makeatother
\subtitle{From Surveillance Capitalism to AI Imperialism \textbar{}
Chapter 06: AI as a dual-use rhetorical weapon}
\author{}
\date{}

\begin{document}
\maketitle

\subsection{Abstract}\label{abstract}

One of the clearest contradictions in the current AI regime is that
systems are marketed simultaneously as dangerous enough to justify elite
governance and safe enough to be rolled out as broad dependence
infrastructure. OpenAI publicly framed reasoning as a distinct
capability layer tied to additional train-time and test-time compute,
and paired it with preparedness, safety, and misuse-risk language. This
is how AI becomes a dual-use rhetorical weapon.

\subsection{Keywords}\label{keywords}

AI as a dual-use rhetorical weapon; surveillance capitalism; AI
imperialism; evidence-based dossier

\subsection{Research Question}\label{research-question}

How does ai as a dual-use rhetorical weapon function within the
dossier's larger argument, and what does the documentary record allow
this chapter to establish?

\subsection{Scope and Framing Note}\label{scope-and-framing-note}

This paper examines ai as a dual-use rhetorical weapon within the
broader dossier \emph{From Surveillance Capitalism to AI Imperialism}.
It is intentionally written to circulate on its own: necessary context
is restated, evidence boundaries are made explicit, and cited material
is reproduced in paper-local form rather than delegated to repo
navigation. It contributes to the series as a paper that shows how
danger rhetoric, cheapness rhetoric, and reasoning rhetoric together
normalize frontier enclosure and strategic dependence. Adjacent
reinforcement appears in Chapter 04, Chapter 05, Chapter 07, Chapter 10,
but those cross-references are supportive rather than required for basic
comprehension.

\subsection{Central Thesis}\label{central-thesis}

One of the clearest contradictions in the current AI regime is that
systems are marketed simultaneously as dangerous enough to justify elite
governance and safe enough to be rolled out as broad dependence
infrastructure. The point is not that reasoning advances are fake. The
point is that real advances are being packaged inside a politics of
fear, prestige, and access management at the same time. \footnote{Sources:
  OpenAI,
  \href{https://openai.com/index/learning-to-reason-with-llms/}{Learning
  to reason with LLMs} (2024-09-12; see
  \hyperref[appendix-a1-learning-to-reason-with-llms]{Appendix A1});
  OpenAI, \href{https://openai.com/index/openai-o1-system-card/}{OpenAI
  o1 System Card} (2024-12-05; see
  \hyperref[appendix-a2-openai-o1-system-card]{Appendix A2}); Financial
  Times,
  \href{https://www.ft.com/content/37ba7236-2a64-4807-b1e1-7e21ee7d0914}{OpenAI
  acknowledges new models increa\ldots{}} (2024-09-13; see
  \hyperref[appendix-a3-openai-acknowledges-new-models-increase-risk-of-misuse-to-create-bioweapons]{Appendix
  A3})}

\subsection{Evidence and Method Note}\label{evidence-and-method-note}

This paper is derived from the chapter corpus for Chapter 06 and is
grounded in the project's structured research OS. It relies on source
notes, extracted claims, entity profiles, and timeline events already
normalized under the dossier's evidence hierarchy. Documented facts are
asserted directly where the record supports them; disputed matters are
marked as such; interpretive claims are bounded explicitly; and
speculative overreach is isolated in the evidence-boundary appendix.

\subsection{Introduction}\label{introduction}

This paper asks: How does ai as a dual-use rhetorical weapon function
within the dossier's larger argument, and what does the documentary
record allow this chapter to establish? Its working answer is clear from
the start: One of the clearest contradictions in the current AI regime
is that systems are marketed simultaneously as dangerous enough to
justify elite governance and safe enough to be rolled out as broad
dependence infrastructure. The point is not that reasoning advances are
fake. The point is that real advances are being packaged inside a
politics of fear, prestige, and access management at the same time.
\footnote{Sources: OpenAI,
  \href{https://openai.com/index/learning-to-reason-with-llms/}{Learning
  to reason with LLMs} (2024-09-12; see
  \hyperref[appendix-a1-learning-to-reason-with-llms]{Appendix A1});
  OpenAI, \href{https://openai.com/index/openai-o1-system-card/}{OpenAI
  o1 System Card} (2024-12-05; see
  \hyperref[appendix-a2-openai-o1-system-card]{Appendix A2}); Financial
  Times,
  \href{https://www.ft.com/content/37ba7236-2a64-4807-b1e1-7e21ee7d0914}{OpenAI
  acknowledges new models increa\ldots{}} (2024-09-13; see
  \hyperref[appendix-a3-openai-acknowledges-new-models-increase-risk-of-misuse-to-create-bioweapons]{Appendix
  A3})} The objective is not only to recount developments, but to show
why the subject of this chapter belongs inside a structural argument
about extraction, enclosure, legitimacy, and reform.

\subsection{Historical or Institutional
Context}\label{historical-or-institutional-context}

OpenAI publicly framed reasoning as a distinct capability layer tied to
additional train-time and test-time compute, and paired it with
preparedness, safety, and misuse-risk language. The product story and
the security story arrived together. \footnote{Sources: OpenAI,
  \href{https://openai.com/index/learning-to-reason-with-llms/}{Learning
  to reason with LLMs} (2024-09-12; see
  \hyperref[appendix-a1-learning-to-reason-with-llms]{Appendix A1});
  OpenAI, \href{https://openai.com/index/openai-o1-system-card/}{OpenAI
  o1 System Card} (2024-12-05; see
  \hyperref[appendix-a2-openai-o1-system-card]{Appendix A2}); Financial
  Times,
  \href{https://www.ft.com/content/37ba7236-2a64-4807-b1e1-7e21ee7d0914}{OpenAI
  acknowledges new models increa\ldots{}} (2024-09-13; see
  \hyperref[appendix-a3-openai-acknowledges-new-models-increase-risk-of-misuse-to-create-bioweapons]{Appendix
  A3})}

\subsection{Main Analysis}\label{main-analysis}

The Reuters Q* and \texttt{Strawberry} reporting shows that this
reasoning story has a serious public prehistory rather than being a
later marketing improvisation. That matters because it confirms both
sides of the contradiction: the capability arc may be real, but so is
the later struggle over how it is narrated, priced, and gated.
\footnote{Sources: Reuters,
  \href{https://www.reuters.com/technology/sam-altmans-ouster-openai-was-precipitated-by-letter-board-about-ai-breakthrough-2023-11-22/}{OpenAI
  researchers warned board of AI\ldots{}} (2023-11-23; see
  \hyperref[appendix-a4-openai-researchers-warned-board-of-ai-breakthrough-ahead-of-ceo-ouster-sources-say]{Appendix
  A4}); Reuters,
  \href{https://www.reuters.com/technology/artificial-intelligence/openai-working-new-reasoning-technology-under-code-name-strawberry-2024-07-12/}{Exclusive:
  OpenAI working on new reas\ldots{}} (2024-07-15; see
  \hyperref[appendix-a5-exclusive-openai-working-on-new-reasoning-technology-under-code-name-strawberry]{Appendix
  A5})}

Primary research strengthens the middle claim. Test-time compute is not
just a corporate slogan. A 2024 paper on compute-optimal inference-time
scaling showed meaningful efficiency improvements over simpler
baselines, while later research and open-replication efforts showed both
rapid diffusion pressure and an active technical dispute over what those
replications actually reproduced. In other words, the dossier no longer
needs to pretend that reasoning scarcity is fake in order to criticize
how it is governed. \footnote{Sources: Charlie Snell; Jaehoon Lee;
  Kelvin Xu; Aviral Kumar,
  \href{https://arxiv.org/abs/2408.03314}{Scaling LLM Test-Time Compute
  Optimal\ldots{}} (2024-08-06; see
  \hyperref[appendix-a6-scaling-llm-test-time-compute-optimally-can-be-more-effective-than-scaling-model-parameters]{Appendix
  A6}); Niklas Muennighoff; Zitong Yang; Weijia Shi; Xiang Lisa Li; Li
  Fei-Fei; Hannaneh Hajishirzi; Luke Zettlemoyer; Percy Liang; Emmanuel
  Candès; Tatsunori Hashimoto,
  \href{https://arxiv.org/abs/2501.19393}{s1: Simple test-time scaling}
  (2025-01-31; see
  \hyperref[appendix-a7-s1-simple-test-time-scaling]{Appendix A7});
  Guojun Wu, \href{https://arxiv.org/abs/2507.14419}{It's Not That
  Simple. An Analysis of\ldots{}} (2025-07-19; see
  \hyperref[appendix-a8-it-s-not-that-simple-an-analysis-of-simple-test-time-scaling]{Appendix
  A8})}

The pricing layer is also direct now. OpenAI's own plan and API pages,
along with Anthropic's pricing and \texttt{Claude\ 4} launch material,
show that stronger reasoning access is tiered across premium plans,
higher usage classes, and higher-priced API lanes. Reasoning is not only
a capability. It is also a gated commercial product. \footnote{Sources:
  OpenAI, \href{https://chatgpt.com/pricing/}{ChatGPT Plans \textbar{}
  Free, Go, Plus, Pro,\ldots{}} (n.d.; see
  \hyperref[appendix-a9-chatgpt-plans-free-go-plus-pro-business-and-enterprise]{Appendix
  A9}); OpenAI,
  \href{https://developers.openai.com/api/docs/pricing}{Pricing
  \textbar{} OpenAI API} (n.d.; see
  \hyperref[appendix-a10-pricing-openai-api]{Appendix A10}); Anthropic,
  \href{https://claude.com/pricing}{Plans \& Pricing \textbar{} Claude
  by Anthropic} (n.d.; see
  \hyperref[appendix-a11-plans-pricing-claude-by-anthropic]{Appendix
  A11}); Anthropic,
  \href{https://www.anthropic.com/news/claude-4}{Introducing Claude 4}
  (2025-05-22; see \hyperref[appendix-a12-introducing-claude-4]{Appendix
  A12})}

DeepSeek complicates total exclusivity claims without eliminating
concentration. Its releases show that parts of reasoning capability
diffuse faster than absolutist frontier rhetoric suggests. The real
story is not zero diffusion versus perfect monopoly. It is diffusion at
the edges and political curation at the core. \footnote{Sources: arXiv,
  \href{https://arxiv.org/abs/2412.19437}{DeepSeek-V3 Technical Report}
  (2024-12-27; see
  \hyperref[appendix-a13-deepseek-v3-technical-report]{Appendix A13});
  arXiv, \href{https://arxiv.org/abs/2501.12948}{DeepSeek-R1:
  Incentivizing Reasoning\ldots{}} (2025-01-22; see
  \hyperref[appendix-a14-deepseek-r1-incentivizing-reasoning-capability-in-llms-via-reinforcement-learning]{Appendix
  A14}); Niklas Muennighoff; Zitong Yang; Weijia Shi; Xiang Lisa Li; Li
  Fei-Fei; Hannaneh Hajishirzi; Luke Zettlemoyer; Percy Liang; Emmanuel
  Candès; Tatsunori Hashimoto,
  \href{https://arxiv.org/abs/2501.19393}{s1: Simple test-time scaling}
  (2025-01-31; see
  \hyperref[appendix-a7-s1-simple-test-time-scaling]{Appendix A7})}

Interpretation: this chapter supports the patrimony argument by showing
how a real compute-intensive technical mechanism is turned into a
managed access regime. Fear legitimizes the gate, convenience fills the
funnel, and premium reasoning tiers convert collectively grounded
capability into selectively governed scarcity. \footnote{Sources:
  OpenAI, \href{https://openai.com/index/openai-o1-system-card/}{OpenAI
  o1 System Card} (2024-12-05; see
  \hyperref[appendix-a2-openai-o1-system-card]{Appendix A2}); OpenAI,
  \href{https://chatgpt.com/pricing/}{ChatGPT Plans \textbar{} Free, Go,
  Plus, Pro,\ldots{}} (n.d.; see
  \hyperref[appendix-a9-chatgpt-plans-free-go-plus-pro-business-and-enterprise]{Appendix
  A9}); Anthropic, \href{https://claude.com/pricing}{Plans \& Pricing
  \textbar{} Claude by Anthropic} (n.d.; see
  \hyperref[appendix-a11-plans-pricing-claude-by-anthropic]{Appendix
  A11}); Charlie Snell; Jaehoon Lee; Kelvin Xu; Aviral Kumar,
  \href{https://arxiv.org/abs/2408.03314}{Scaling LLM Test-Time Compute
  Optimal\ldots{}} (2024-08-06; see
  \hyperref[appendix-a6-scaling-llm-test-time-compute-optimally-can-be-more-effective-than-scaling-model-parameters]{Appendix
  A6})}

\subsection{Position Within the
Series}\label{position-within-the-series}

This paper sits inside a coordinated 11-paper architecture. Its nearest
companions are Chapter 04 (AI arrives when the human raw material
already exists), Chapter 05 (Talent, overhiring, valuation, and
layoffs), Chapter 07 (Imperial control over access), Chapter 10 (What to
do). In that series logic, Chapter 06 shows how danger rhetoric,
cheapness rhetoric, and reasoning rhetoric together normalize frontier
enclosure and strategic dependence.

\subsection{Counterarguments}\label{counterarguments}

The chapter should not collapse three separate issues into one: the
reality of test-time compute as a technical mechanism, the open
challenge to frontier mystique, and the unresolved dispute over what
simple replications actually prove. The strongest claim is about
rhetoric, gating, and governance, not about proving that every frontier
capability turned out to be hollow. \footnote{Sources: Charlie Snell;
  Jaehoon Lee; Kelvin Xu; Aviral Kumar,
  \href{https://arxiv.org/abs/2408.03314}{Scaling LLM Test-Time Compute
  Optimal\ldots{}} (2024-08-06; see
  \hyperref[appendix-a6-scaling-llm-test-time-compute-optimally-can-be-more-effective-than-scaling-model-parameters]{Appendix
  A6}); Niklas Muennighoff; Zitong Yang; Weijia Shi; Xiang Lisa Li; Li
  Fei-Fei; Hannaneh Hajishirzi; Luke Zettlemoyer; Percy Liang; Emmanuel
  Candès; Tatsunori Hashimoto,
  \href{https://arxiv.org/abs/2501.19393}{s1: Simple test-time scaling}
  (2025-01-31; see
  \hyperref[appendix-a7-s1-simple-test-time-scaling]{Appendix A7});
  Guojun Wu, \href{https://arxiv.org/abs/2507.14419}{It's Not That
  Simple. An Analysis of\ldots{}} (2025-07-19; see
  \hyperref[appendix-a8-it-s-not-that-simple-an-analysis-of-simple-test-time-scaling]{Appendix
  A8})}

\subsection{Limits of the Evidence}\label{limits-of-the-evidence}

The chapter should not collapse three separate issues into one: the
reality of test-time compute as a technical mechanism, the open
challenge to frontier mystique, and the unresolved dispute over what
simple replications actually prove. The strongest claim is about
rhetoric, gating, and governance, not about proving that every frontier
capability turned out to be hollow. \footnote{Sources: Charlie Snell;
  Jaehoon Lee; Kelvin Xu; Aviral Kumar,
  \href{https://arxiv.org/abs/2408.03314}{Scaling LLM Test-Time Compute
  Optimal\ldots{}} (2024-08-06; see
  \hyperref[appendix-a6-scaling-llm-test-time-compute-optimally-can-be-more-effective-than-scaling-model-parameters]{Appendix
  A6}); Niklas Muennighoff; Zitong Yang; Weijia Shi; Xiang Lisa Li; Li
  Fei-Fei; Hannaneh Hajishirzi; Luke Zettlemoyer; Percy Liang; Emmanuel
  Candès; Tatsunori Hashimoto,
  \href{https://arxiv.org/abs/2501.19393}{s1: Simple test-time scaling}
  (2025-01-31; see
  \hyperref[appendix-a7-s1-simple-test-time-scaling]{Appendix A7});
  Guojun Wu, \href{https://arxiv.org/abs/2507.14419}{It's Not That
  Simple. An Analysis of\ldots{}} (2025-07-19; see
  \hyperref[appendix-a8-it-s-not-that-simple-an-analysis-of-simple-test-time-scaling]{Appendix
  A8})}

\subsection{Reform Relevance}\label{reform-relevance}

This is how AI becomes a dual-use rhetorical weapon. Danger justifies
concentration, while convenience justifies mass adoption. A public told
that the system is too dangerous to democratize and too useful to refuse
is being prepared to accept hierarchy as prudence. That pushes Chapter
\texttt{04}'s ownership question toward a sharper answer: capability
built from collective human production is being enclosed as a governed
premium. \footnote{Sources: OpenAI,
  \href{https://openai.com/index/openai-o1-system-card/}{OpenAI o1
  System Card} (2024-12-05; see
  \hyperref[appendix-a2-openai-o1-system-card]{Appendix A2}); OpenAI,
  \href{https://chatgpt.com/pricing/}{ChatGPT Plans \textbar{} Free, Go,
  Plus, Pro,\ldots{}} (n.d.; see
  \hyperref[appendix-a9-chatgpt-plans-free-go-plus-pro-business-and-enterprise]{Appendix
  A9}); Anthropic,
  \href{https://www.anthropic.com/news/claude-4}{Introducing Claude 4}
  (2025-05-22; see \hyperref[appendix-a12-introducing-claude-4]{Appendix
  A12})}

\subsection{Conclusion}\label{conclusion}

This is how AI becomes a dual-use rhetorical weapon. Once managed
dependency exists, the next question is who controls the most
consequential access lanes and under what geopolitical terms. In the
architecture of this series, Chapter 06 therefore functions as a
self-contained argument while also advancing the cumulative path toward
the dossier's three reforms.

\subsection{Notes}\label{notes}

This paper is designed to circulate independently. Public notes point
readers to the real external documents first, then to the paper-local
appendix entry that explains why each source matters.

The underlying research OS distinguishes among
\texttt{Documented\ Fact}, \texttt{Disputed\ Fact},
\texttt{Hypothesis\ /\ Interpretation}, and
\texttt{Speculative\ Narrative\ Risk}. In Chapter 06, those boundaries
remain visible in the prose, in the bibliography, and in the appendix
sections that summarize claims and caution points.

\subsection{References / Bibliography}\label{references-bibliography}

\begin{enumerate}
\def\labelenumi{\arabic{enumi}.}
\tightlist
\item
  OpenAI.
  \href{https://openai.com/index/learning-to-reason-with-llms/}{Learning
  to reason with LLMs}. 2024-09-12. T1 release. Paper appendix:
  \hyperref[appendix-a1-learning-to-reason-with-llms]{Appendix A1}.
  Corpus companion entry: Source S-093. Audit ID:
  \texttt{openai-learning-to-reason-with-llms-2024}. Relevance: That
  OpenAI publicly framed reasoning as a new scaling layer based on
  additional train-time and test-time compute.
\item
  OpenAI. \href{https://openai.com/index/openai-o1-system-card/}{OpenAI
  o1 System Card}. 2024-12-05. T1 system\_card. Paper appendix:
  \hyperref[appendix-a2-openai-o1-system-card]{Appendix A2}. Corpus
  companion entry: Source S-095. Audit ID:
  \texttt{openai-o1-system-card-2024}. Relevance: That OpenAI itself
  paired stronger reasoning capability with formal safety, preparedness,
  and misuse-risk framing.
\item
  Financial Times.
  \href{https://www.ft.com/content/37ba7236-2a64-4807-b1e1-7e21ee7d0914}{OpenAI
  acknowledges new models increase risk of misuse to create bioweapons}.
  2024-09-13. T2 news\_article. Paper appendix:
  \hyperref[appendix-a3-openai-acknowledges-new-models-increase-risk-of-misuse-to-create-bioweapons]{Appendix
  A3}. Corpus companion entry: Source S-060. Audit ID:
  \texttt{ft-openai-o1-bioweapon-risk-2024}. Relevance: That serious
  \texttt{T2} reporting documented OpenAI publicly acknowledging higher
  misuse risk at the same moment it was releasing stronger reasoning
  models.
\item
  Reuters.
  \href{https://www.reuters.com/technology/sam-altmans-ouster-openai-was-precipitated-by-letter-board-about-ai-breakthrough-2023-11-22/}{OpenAI
  researchers warned board of AI breakthrough ahead of CEO ouster,
  sources say}. 2023-11-23. T2 news\_article. Paper appendix:
  \hyperref[appendix-a4-openai-researchers-warned-board-of-ai-breakthrough-ahead-of-ceo-ouster-sources-say]{Appendix
  A4}. Corpus companion entry: Source S-110. Audit ID:
  \texttt{reuters-qstar-board-warning-2023}. Relevance: That by November
  23, 2023 there was serious mainstream reporting explicitly linking the
  board crisis to a letter about an AI breakthrough.
\item
  Reuters.
  \href{https://www.reuters.com/technology/artificial-intelligence/openai-working-new-reasoning-technology-under-code-name-strawberry-2024-07-12/}{Exclusive:
  OpenAI working on new reasoning technology under code name
  `Strawberry'}. 2024-07-15. T2 news\_article. Paper appendix:
  \hyperref[appendix-a5-exclusive-openai-working-on-new-reasoning-technology-under-code-name-strawberry]{Appendix
  A5}. Corpus companion entry: Source S-111. Audit ID:
  \texttt{reuters-strawberry-reasoning-2024}. Relevance: That by
  mid-July 2024 Reuters had published serious reporting about an OpenAI
  reasoning project under the codename \texttt{Strawberry}.
\item
  Charlie Snell; Jaehoon Lee; Kelvin Xu; Aviral Kumar.
  \href{https://arxiv.org/abs/2408.03314}{Scaling LLM Test-Time Compute
  Optimally can be More Effective than Scaling Model Parameters}.
  2024-08-06. T1 paper. Paper appendix:
  \hyperref[appendix-a6-scaling-llm-test-time-compute-optimally-can-be-more-effective-than-scaling-model-parameters]{Appendix
  A6}. Corpus companion entry: Source S-115. Audit ID:
  \texttt{snell-test-time-compute-2024}. Relevance: That test-time
  compute scaling is a documented technical mechanism rather than only a
  commercial branding story.
\item
  Niklas Muennighoff; Zitong Yang; Weijia Shi; Xiang Lisa Li; Li
  Fei-Fei; Hannaneh Hajishirzi; Luke Zettlemoyer; Percy Liang; Emmanuel
  Candès; Tatsunori Hashimoto.
  \href{https://arxiv.org/abs/2501.19393}{s1: Simple test-time scaling}.
  2025-01-31. T1 paper. Paper appendix:
  \hyperref[appendix-a7-s1-simple-test-time-scaling]{Appendix A7}.
  Corpus companion entry: Source S-112. Audit ID:
  \texttt{s1-simple-test-time-scaling-2025}. Relevance: That OpenAI's
  reasoning launch quickly triggered open replication efforts focused on
  test-time scaling.
\item
  Guojun Wu. \href{https://arxiv.org/abs/2507.14419}{It's Not That
  Simple. An Analysis of Simple Test-Time Scaling}. 2025-07-19. T1
  paper. Paper appendix:
  \hyperref[appendix-a8-it-s-not-that-simple-an-analysis-of-simple-test-time-scaling]{Appendix
  A8}. Corpus companion entry: Source S-076. Audit ID:
  \texttt{its-not-that-simple-test-time-scaling-2025}. Relevance: That
  there is direct technical pushback against simplistic claims that
  frontier reasoning was trivially replicated.
\item
  OpenAI. \href{https://chatgpt.com/pricing/}{ChatGPT Plans \textbar{}
  Free, Go, Plus, Pro, Business, and Enterprise}. n.d.. T1
  pricing\_page. Paper appendix:
  \hyperref[appendix-a9-chatgpt-plans-free-go-plus-pro-business-and-enterprise]{Appendix
  A9}. Corpus companion entry: Source S-039. Audit ID:
  \texttt{chatgpt-pricing-2026}. Relevance: That as of capture on July
  3, 2026, OpenAI's official ChatGPT plan structure reserved the
  highest-end reasoning access for upper paid tiers rather than offering
  it uniformly acr\ldots{}
\item
  OpenAI. \href{https://developers.openai.com/api/docs/pricing}{Pricing
  \textbar{} OpenAI API}. n.d.. T1 pricing\_page. Paper appendix:
  \hyperref[appendix-a10-pricing-openai-api]{Appendix A10}. Corpus
  companion entry: Source S-088. Audit ID:
  \texttt{openai-api-pricing-2026}. Relevance: That as of July 3, 2026,
  OpenAI's official API pricing imposed a large price premium on
  \texttt{GPT-5.5\ Pro} relative to \texttt{GPT-5.5}.
\item
  Anthropic. \href{https://claude.com/pricing}{Plans \& Pricing
  \textbar{} Claude by Anthropic}. n.d.. T1 pricing\_page. Paper
  appendix:
  \hyperref[appendix-a11-plans-pricing-claude-by-anthropic]{Appendix
  A11}. Corpus companion entry: Source S-005. Audit ID:
  \texttt{anthropic-claude-pricing-2026}. Relevance: That as of July 3,
  2026, Anthropic's official current pricing structure tied higher
  usage, priority access, and stronger feature access to premium paid
  plans.
\item
  Anthropic. \href{https://www.anthropic.com/news/claude-4}{Introducing
  Claude 4}. 2025-05-22. T1 release. Paper appendix:
  \hyperref[appendix-a12-introducing-claude-4]{Appendix A12}. Corpus
  companion entry: Source S-003. Audit ID:
  \texttt{anthropic-claude-4-release-2025}. Relevance: That on May 22,
  2025 Anthropic launched reasoning-oriented hybrid frontier models
  together with immediate commercial distribution across subscription
  plans, API access, and majo\ldots{}
\item
  arXiv. \href{https://arxiv.org/abs/2412.19437}{DeepSeek-V3 Technical
  Report}. 2024-12-27. T1 paper. Paper appendix:
  \hyperref[appendix-a13-deepseek-v3-technical-report]{Appendix A13}.
  Corpus companion entry: Source S-049. Audit ID:
  \texttt{deepseek-v3-technical-report-2024}. Relevance: That by late
  2024 there was already a competitor publicly claiming
  frontier-adjacent performance through a technical report.
\item
  arXiv. \href{https://arxiv.org/abs/2501.12948}{DeepSeek-R1:
  Incentivizing Reasoning Capability in LLMs via Reinforcement
  Learning}. 2025-01-22. T1 paper. Paper appendix:
  \hyperref[appendix-a14-deepseek-r1-incentivizing-reasoning-capability-in-llms-via-reinforcement-learning]{Appendix
  A14}. Corpus companion entry: Source S-048. Audit ID:
  \texttt{deepseek-r1-reinforcement-learning-2025}. Relevance: That by
  January 2025 a non-U.S. lab was publicly claiming frontier-adjacent
  reasoning performance in a primary technical paper.
\end{enumerate}

\subsection{Appendix A. Source Register for This
Paper}\label{appendix-a.-source-register-for-this-paper}

\subsubsection{Appendix A1. Learning to reason with
LLMs}\label{appendix-a1-learning-to-reason-with-llms}

\textbf{Source citation:}
\href{https://openai.com/index/learning-to-reason-with-llms/}{Learning
to reason with LLMs} \textbf{Institution / author:} OpenAI
\textbf{Date:} 2024-09-12 \textbf{Tier / type:} T1 / release \textbf{Why
it matters in this paper:} That OpenAI publicly framed reasoning as a
new scaling layer based on additional train-time and test-time compute.
\textbf{Corpus companion entry:} Source S-093 \textbf{Audit ID:}
\texttt{openai-learning-to-reason-with-llms-2024}

\subsubsection{Appendix A2. OpenAI o1 System
Card}\label{appendix-a2-openai-o1-system-card}

\textbf{Source citation:}
\href{https://openai.com/index/openai-o1-system-card/}{OpenAI o1 System
Card} \textbf{Institution / author:} OpenAI \textbf{Date:} 2024-12-05
\textbf{Tier / type:} T1 / system\_card \textbf{Why it matters in this
paper:} That OpenAI itself paired stronger reasoning capability with
formal safety, preparedness, and misuse-risk framing. \textbf{Corpus
companion entry:} Source S-095 \textbf{Audit ID:}
\texttt{openai-o1-system-card-2024}

\subsubsection{Appendix A3. OpenAI acknowledges new models increase risk
of misuse to create
bioweapons}\label{appendix-a3-openai-acknowledges-new-models-increase-risk-of-misuse-to-create-bioweapons}

\textbf{Source citation:}
\href{https://www.ft.com/content/37ba7236-2a64-4807-b1e1-7e21ee7d0914}{OpenAI
acknowledges new models increase risk of misuse to create bioweapons}
\textbf{Institution / author:} Financial Times \textbf{Date:} 2024-09-13
\textbf{Tier / type:} T2 / news\_article \textbf{Why it matters in this
paper:} That serious \texttt{T2} reporting documented OpenAI publicly
acknowledging higher misuse risk at the same moment it was releasing
stronger reasoning models. \textbf{Corpus companion entry:} Source S-060
\textbf{Audit ID:} \texttt{ft-openai-o1-bioweapon-risk-2024}

\subsubsection{Appendix A4. OpenAI researchers warned board of AI
breakthrough ahead of CEO ouster, sources
say}\label{appendix-a4-openai-researchers-warned-board-of-ai-breakthrough-ahead-of-ceo-ouster-sources-say}

\textbf{Source citation:}
\href{https://www.reuters.com/technology/sam-altmans-ouster-openai-was-precipitated-by-letter-board-about-ai-breakthrough-2023-11-22/}{OpenAI
researchers warned board of AI breakthrough ahead of CEO ouster, sources
say} \textbf{Institution / author:} Reuters \textbf{Date:} 2023-11-23
\textbf{Tier / type:} T2 / news\_article \textbf{Why it matters in this
paper:} That by November 23, 2023 there was serious mainstream reporting
explicitly linking the board crisis to a letter about an AI
breakthrough. \textbf{Corpus companion entry:} Source S-110
\textbf{Audit ID:} \texttt{reuters-qstar-board-warning-2023}

\subsubsection{Appendix A5. Exclusive: OpenAI working on new reasoning
technology under code name
`Strawberry'}\label{appendix-a5-exclusive-openai-working-on-new-reasoning-technology-under-code-name-strawberry}

\textbf{Source citation:}
\href{https://www.reuters.com/technology/artificial-intelligence/openai-working-new-reasoning-technology-under-code-name-strawberry-2024-07-12/}{Exclusive:
OpenAI working on new reasoning technology under code name `Strawberry'}
\textbf{Institution / author:} Reuters \textbf{Date:} 2024-07-15
\textbf{Tier / type:} T2 / news\_article \textbf{Why it matters in this
paper:} That by mid-July 2024 Reuters had published serious reporting
about an OpenAI reasoning project under the codename
\texttt{Strawberry}. \textbf{Corpus companion entry:} Source S-111
\textbf{Audit ID:} \texttt{reuters-strawberry-reasoning-2024}

\subsubsection{Appendix A6. Scaling LLM Test-Time Compute Optimally can
be More Effective than Scaling Model
Parameters}\label{appendix-a6-scaling-llm-test-time-compute-optimally-can-be-more-effective-than-scaling-model-parameters}

\textbf{Source citation:}
\href{https://arxiv.org/abs/2408.03314}{Scaling LLM Test-Time Compute
Optimally can be More Effective than Scaling Model Parameters}
\textbf{Institution / author:} Charlie Snell; Jaehoon Lee; Kelvin Xu;
Aviral Kumar \textbf{Date:} 2024-08-06 \textbf{Tier / type:} T1 / paper
\textbf{Why it matters in this paper:} That test-time compute scaling is
a documented technical mechanism rather than only a commercial branding
story. \textbf{Corpus companion entry:} Source S-115 \textbf{Audit ID:}
\texttt{snell-test-time-compute-2024}

\subsubsection{Appendix A7. s1: Simple test-time
scaling}\label{appendix-a7-s1-simple-test-time-scaling}

\textbf{Source citation:} \href{https://arxiv.org/abs/2501.19393}{s1:
Simple test-time scaling} \textbf{Institution / author:} Niklas
Muennighoff; Zitong Yang; Weijia Shi; Xiang Lisa Li; Li Fei-Fei;
Hannaneh Hajishirzi; Luke Zettlemoyer; Percy Liang; Emmanuel Candès;
Tatsunori Hashimoto \textbf{Date:} 2025-01-31 \textbf{Tier / type:} T1 /
paper \textbf{Why it matters in this paper:} That OpenAI's reasoning
launch quickly triggered open replication efforts focused on test-time
scaling. \textbf{Corpus companion entry:} Source S-112 \textbf{Audit
ID:} \texttt{s1-simple-test-time-scaling-2025}

\subsubsection{Appendix A8. It's Not That Simple. An Analysis of Simple
Test-Time
Scaling}\label{appendix-a8-it-s-not-that-simple-an-analysis-of-simple-test-time-scaling}

\textbf{Source citation:} \href{https://arxiv.org/abs/2507.14419}{It's
Not That Simple. An Analysis of Simple Test-Time Scaling}
\textbf{Institution / author:} Guojun Wu \textbf{Date:} 2025-07-19
\textbf{Tier / type:} T1 / paper \textbf{Why it matters in this paper:}
That there is direct technical pushback against simplistic claims that
frontier reasoning was trivially replicated. \textbf{Corpus companion
entry:} Source S-076 \textbf{Audit ID:}
\texttt{its-not-that-simple-test-time-scaling-2025}

\subsubsection{Appendix A9. ChatGPT Plans \textbar{} Free, Go, Plus,
Pro, Business, and
Enterprise}\label{appendix-a9-chatgpt-plans-free-go-plus-pro-business-and-enterprise}

\textbf{Source citation:} \href{https://chatgpt.com/pricing/}{ChatGPT
Plans \textbar{} Free, Go, Plus, Pro, Business, and Enterprise}
\textbf{Institution / author:} OpenAI \textbf{Date:} n.d. \textbf{Tier /
type:} T1 / pricing\_page \textbf{Why it matters in this paper:} That as
of capture on July 3, 2026, OpenAI's official ChatGPT plan structure
reserved the highest-end reasoning access for upper paid tiers rather
than offering it uniformly acr\ldots{} \textbf{Corpus companion entry:}
Source S-039 \textbf{Audit ID:} \texttt{chatgpt-pricing-2026}

\subsubsection{Appendix A10. Pricing \textbar{} OpenAI
API}\label{appendix-a10-pricing-openai-api}

\textbf{Source citation:}
\href{https://developers.openai.com/api/docs/pricing}{Pricing \textbar{}
OpenAI API} \textbf{Institution / author:} OpenAI \textbf{Date:} n.d.
\textbf{Tier / type:} T1 / pricing\_page \textbf{Why it matters in this
paper:} That as of July 3, 2026, OpenAI's official API pricing imposed a
large price premium on \texttt{GPT-5.5\ Pro} relative to
\texttt{GPT-5.5}. \textbf{Corpus companion entry:} Source S-088
\textbf{Audit ID:} \texttt{openai-api-pricing-2026}

\subsubsection{Appendix A11. Plans \& Pricing \textbar{} Claude by
Anthropic}\label{appendix-a11-plans-pricing-claude-by-anthropic}

\textbf{Source citation:} \href{https://claude.com/pricing}{Plans \&
Pricing \textbar{} Claude by Anthropic} \textbf{Institution / author:}
Anthropic \textbf{Date:} n.d. \textbf{Tier / type:} T1 / pricing\_page
\textbf{Why it matters in this paper:} That as of July 3, 2026,
Anthropic's official current pricing structure tied higher usage,
priority access, and stronger feature access to premium paid plans.
\textbf{Corpus companion entry:} Source S-005 \textbf{Audit ID:}
\texttt{anthropic-claude-pricing-2026}

\subsubsection{Appendix A12. Introducing Claude
4}\label{appendix-a12-introducing-claude-4}

\textbf{Source citation:}
\href{https://www.anthropic.com/news/claude-4}{Introducing Claude 4}
\textbf{Institution / author:} Anthropic \textbf{Date:} 2025-05-22
\textbf{Tier / type:} T1 / release \textbf{Why it matters in this
paper:} That on May 22, 2025 Anthropic launched reasoning-oriented
hybrid frontier models together with immediate commercial distribution
across subscription plans, API access, and majo\ldots{} \textbf{Corpus
companion entry:} Source S-003 \textbf{Audit ID:}
\texttt{anthropic-claude-4-release-2025}

\subsubsection{Appendix A13. DeepSeek-V3 Technical
Report}\label{appendix-a13-deepseek-v3-technical-report}

\textbf{Source citation:}
\href{https://arxiv.org/abs/2412.19437}{DeepSeek-V3 Technical Report}
\textbf{Institution / author:} arXiv \textbf{Date:} 2024-12-27
\textbf{Tier / type:} T1 / paper \textbf{Why it matters in this paper:}
That by late 2024 there was already a competitor publicly claiming
frontier-adjacent performance through a technical report. \textbf{Corpus
companion entry:} Source S-049 \textbf{Audit ID:}
\texttt{deepseek-v3-technical-report-2024}

\subsubsection{Appendix A14. DeepSeek-R1: Incentivizing Reasoning
Capability in LLMs via Reinforcement
Learning}\label{appendix-a14-deepseek-r1-incentivizing-reasoning-capability-in-llms-via-reinforcement-learning}

\textbf{Source citation:}
\href{https://arxiv.org/abs/2501.12948}{DeepSeek-R1: Incentivizing
Reasoning Capability in LLMs via Reinforcement Learning}
\textbf{Institution / author:} arXiv \textbf{Date:} 2025-01-22
\textbf{Tier / type:} T1 / paper \textbf{Why it matters in this paper:}
That by January 2025 a non-U.S. lab was publicly claiming
frontier-adjacent reasoning performance in a primary technical paper.
\textbf{Corpus companion entry:} Source S-048 \textbf{Audit ID:}
\texttt{deepseek-r1-reinforcement-learning-2025}

\subsection{Appendix B. Claims Used in This
Paper}\label{appendix-b.-claims-used-in-this-paper}

\subsubsection{Appendix B1. On May 22, 2025 Anthropic launched
reasoning-oriented hybrid frontier models
toge\ldots{}}\label{appendix-b1-on-may-22-2025-anthropic-launched-reasoning-oriented-hybrid-frontier-models-toge}

\textbf{Claim statement:} That on May 22, 2025 Anthropic launched
reasoning-oriented hybrid frontier models together with immediate
commercial distribution across subscription plans, API access, and major
cloud platforms. \textbf{Evidence status:} Documented Fact
\textbf{Source support:}
\hyperref[appendix-a12-introducing-claude-4]{Appendix A12}
\textbf{Corpus companion entry:} Claim C-007 \textbf{Audit Claim ID:}
\texttt{anthropic-claude-4-release-2025-06-ai-double-use-reasoning-deepseek-claim-01}

\subsubsection{Appendix B2. Anthropic itself tied extended-thinking
capability to differentiated plan
availab\ldots{}}\label{appendix-b2-anthropic-itself-tied-extended-thinking-capability-to-differentiated-plan-availab}

\textbf{Claim statement:} That Anthropic itself tied extended-thinking
capability to differentiated plan availability, with full-model access
concentrated in paid tiers even while Sonnet 4 also reached free users.
\textbf{Evidence status:} Documented Fact \textbf{Source support:}
\hyperref[appendix-a12-introducing-claude-4]{Appendix A12}
\textbf{Corpus companion entry:} Claim C-008 \textbf{Audit Claim ID:}
\texttt{anthropic-claude-4-release-2025-06-ai-double-use-reasoning-deepseek-claim-02}

\subsubsection{Appendix B3. Reasoning-capable frontier models were
commercialized with explicit API price
dif\ldots{}}\label{appendix-b3-reasoning-capable-frontier-models-were-commercialized-with-explicit-api-price-dif}

\textbf{Claim statement:} That reasoning-capable frontier models were
commercialized with explicit API price differentiation at launch rather
than released only as a research artifact. \textbf{Evidence status:}
Documented Fact \textbf{Source support:}
\hyperref[appendix-a12-introducing-claude-4]{Appendix A12}
\textbf{Corpus companion entry:} Claim C-009 \textbf{Audit Claim ID:}
\texttt{anthropic-claude-4-release-2025-06-ai-double-use-reasoning-deepseek-claim-03}

\subsubsection{Appendix B4. As of July 3, 2026, Anthropic's official
current pricing structure tied higher
us\ldots{}}\label{appendix-b4-as-of-july-3-2026-anthropic-s-official-current-pricing-structure-tied-higher-us}

\textbf{Claim statement:} That as of July 3, 2026, Anthropic's official
current pricing structure tied higher usage, priority access, and
stronger feature access to premium paid plans. \textbf{Evidence status:}
Documented Fact \textbf{Source support:}
\hyperref[appendix-a11-plans-pricing-claude-by-anthropic]{Appendix A11}
\textbf{Corpus companion entry:} Claim C-013 \textbf{Audit Claim ID:}
\texttt{anthropic-claude-pricing-2026-06-ai-double-use-reasoning-deepseek-claim-01}

\subsubsection{Appendix B5. Anthropic's own current pricing also
differentiates materially across model
famil\ldots{}}\label{appendix-b5-anthropic-s-own-current-pricing-also-differentiates-materially-across-model-famil}

\textbf{Claim statement:} That Anthropic's own current pricing also
differentiates materially across model families at the API level rather
than presenting frontier access as uniform. \textbf{Evidence status:}
Documented Fact \textbf{Source support:}
\hyperref[appendix-a11-plans-pricing-claude-by-anthropic]{Appendix A11}
\textbf{Corpus companion entry:} Claim C-014 \textbf{Audit Claim ID:}
\texttt{anthropic-claude-pricing-2026-06-ai-double-use-reasoning-deepseek-claim-02}

\subsubsection{Appendix B6. Chapter 06 can compare more than one
frontier lab when arguing that
reasoning-era\ldots{}}\label{appendix-b6-chapter-06-can-compare-more-than-one-frontier-lab-when-arguing-that-reasoning-era}

\textbf{Claim statement:} That Chapter 06 can compare more than one
frontier lab when arguing that reasoning-era capability is stratified by
availability and price. \textbf{Evidence status:} Documented Fact
\textbf{Source support:}
\hyperref[appendix-a11-plans-pricing-claude-by-anthropic]{Appendix A11}
\textbf{Corpus companion entry:} Claim C-015 \textbf{Audit Claim ID:}
\texttt{anthropic-claude-pricing-2026-06-ai-double-use-reasoning-deepseek-claim-03}

\subsubsection{Appendix B7. As of capture on July 3, 2026, OpenAI's
official ChatGPT plan structure
reserved\ldots{}}\label{appendix-b7-as-of-capture-on-july-3-2026-openai-s-official-chatgpt-plan-structure-reserved}

\textbf{Claim statement:} That as of capture on July 3, 2026, OpenAI's
official ChatGPT plan structure reserved the highest-end reasoning
access for upper paid tiers rather than offering it uniformly across all
users. \textbf{Evidence status:} Documented Fact \textbf{Source
support:}
\hyperref[appendix-a9-chatgpt-plans-free-go-plus-pro-business-and-enterprise]{Appendix
A9} \textbf{Corpus companion entry:} Claim C-117 \textbf{Audit Claim
ID:}
\texttt{chatgpt-pricing-2026-06-ai-double-use-reasoning-deepseek-claim-01}

\subsubsection{Appendix B8. Current official consumer access to
reasoning-capable models is stratified by
pla\ldots{}}\label{appendix-b8-current-official-consumer-access-to-reasoning-capable-models-is-stratified-by-pla}

\textbf{Claim statement:} That current official consumer access to
reasoning-capable models is stratified by plan, usage envelope, and
feature set. \textbf{Evidence status:} Documented Fact \textbf{Source
support:}
\hyperref[appendix-a9-chatgpt-plans-free-go-plus-pro-business-and-enterprise]{Appendix
A9} \textbf{Corpus companion entry:} Claim C-118 \textbf{Audit Claim
ID:}
\texttt{chatgpt-pricing-2026-06-ai-double-use-reasoning-deepseek-claim-02}

\subsubsection{Appendix B9. Chapter 06 can ground its reasoning-premium
argument in a primary access-policy
s\ldots{}}\label{appendix-b9-chapter-06-can-ground-its-reasoning-premium-argument-in-a-primary-access-policy-s}

\textbf{Claim statement:} That Chapter 06 can ground its
reasoning-premium argument in a primary access-policy source rather than
relying only on commentary about hype or control. \textbf{Evidence
status:} Documented Fact \textbf{Source support:}
\hyperref[appendix-a9-chatgpt-plans-free-go-plus-pro-business-and-enterprise]{Appendix
A9} \textbf{Corpus companion entry:} Claim C-119 \textbf{Audit Claim
ID:}
\texttt{chatgpt-pricing-2026-06-ai-double-use-reasoning-deepseek-claim-03}

\subsubsection{Appendix B10. By January 2025 a non-U.S. lab was publicly
claiming frontier-adjacent
reasoning\ldots{}}\label{appendix-b10-by-january-2025-a-non-u-s-lab-was-publicly-claiming-frontier-adjacent-reasoning}

\textbf{Claim statement:} That by January 2025 a non-U.S. lab was
publicly claiming frontier-adjacent reasoning performance in a primary
technical paper. \textbf{Evidence status:} Documented Fact
\textbf{Source support:}
\hyperref[appendix-a14-deepseek-r1-incentivizing-reasoning-capability-in-llms-via-reinforcement-learning]{Appendix
A14} \textbf{Corpus companion entry:} Claim C-146 \textbf{Audit Claim
ID:}
\texttt{deepseek-r1-reinforcement-learning-2025-06-ai-double-use-reasoning-deepseek-claim-01}

\subsubsection{Appendix B11. Reasoning capability was not confined
entirely to closed U.S. product stacks
and\ldots{}}\label{appendix-b11-reasoning-capability-was-not-confined-entirely-to-closed-u-s-product-stacks-and}

\textbf{Claim statement:} That reasoning capability was not confined
entirely to closed U.S. product stacks and was partly diffusing through
open publication and model release. \textbf{Evidence status:} Documented
Fact \textbf{Source support:}
\hyperref[appendix-a14-deepseek-r1-incentivizing-reasoning-capability-in-llms-via-reinforcement-learning]{Appendix
A14} \textbf{Corpus companion entry:} Claim C-147 \textbf{Audit Claim
ID:}
\texttt{deepseek-r1-reinforcement-learning-2025-06-ai-double-use-reasoning-deepseek-claim-02}

\subsubsection{Appendix B12. Chapters 06 and 07 can contrast
frontier-access politics with a documented
patter\ldots{}}\label{appendix-b12-chapters-06-and-07-can-contrast-frontier-access-politics-with-a-documented-patter}

\textbf{Claim statement:} That Chapters 06 and 07 can contrast
frontier-access politics with a documented pattern of partial technical
diffusion. \textbf{Evidence status:} Documented Fact \textbf{Source
support:}
\hyperref[appendix-a14-deepseek-r1-incentivizing-reasoning-capability-in-llms-via-reinforcement-learning]{Appendix
A14} \textbf{Corpus companion entry:} Claim C-148 \textbf{Audit Claim
ID:}
\texttt{deepseek-r1-reinforcement-learning-2025-06-ai-double-use-reasoning-deepseek-claim-03}

\subsubsection{Appendix B13. By late 2024 there was already a competitor
publicly claiming frontier-adjacent
p\ldots{}}\label{appendix-b13-by-late-2024-there-was-already-a-competitor-publicly-claiming-frontier-adjacent-p}

\textbf{Claim statement:} That by late 2024 there was already a
competitor publicly claiming frontier-adjacent performance through a
technical report. \textbf{Evidence status:} Documented Fact
\textbf{Source support:}
\hyperref[appendix-a13-deepseek-v3-technical-report]{Appendix A13}
\textbf{Corpus companion entry:} Claim C-149 \textbf{Audit Claim ID:}
\texttt{deepseek-v3-technical-report-2024-claim-01}

\subsubsection{Appendix B14. The debate over reasoning, efficiency, and
cost cannot be reduced to the
marketin\ldots{}}\label{appendix-b14-the-debate-over-reasoning-efficiency-and-cost-cannot-be-reduced-to-the-marketin}

\textbf{Claim statement:} That the debate over reasoning, efficiency,
and cost cannot be reduced to the marketing of a single Western actor.
\textbf{Evidence status:} Documented Fact \textbf{Source support:}
\hyperref[appendix-a13-deepseek-v3-technical-report]{Appendix A13}
\textbf{Corpus companion entry:} Claim C-150 \textbf{Audit Claim ID:}
\texttt{deepseek-v3-technical-report-2024-claim-02}

\subsubsection{Appendix B15. Narratives of absolute scarcity around
advanced capabilities must be contrasted
w\ldots{}}\label{appendix-b15-narratives-of-absolute-scarcity-around-advanced-capabilities-must-be-contrasted-w}

\textbf{Claim statement:} That narratives of absolute scarcity around
advanced capabilities must be contrasted with open technical reports and
partial replicability. \textbf{Evidence status:} Documented Fact
\textbf{Source support:}
\hyperref[appendix-a13-deepseek-v3-technical-report]{Appendix A13}
\textbf{Corpus companion entry:} Claim C-151 \textbf{Audit Claim ID:}
\texttt{deepseek-v3-technical-report-2024-claim-03}

\subsubsection{Appendix B16. Serious T2 reporting documented OpenAI
publicly acknowledging higher misuse
risk\ldots{}}\label{appendix-b16-serious-t2-reporting-documented-openai-publicly-acknowledging-higher-misuse-risk}

\textbf{Claim statement:} That serious \texttt{T2} reporting documented
OpenAI publicly acknowledging higher misuse risk at the same moment it
was releasing stronger reasoning models. \textbf{Evidence status:}
Documented Fact \textbf{Source support:}
\hyperref[appendix-a3-openai-acknowledges-new-models-increase-risk-of-misuse-to-create-bioweapons]{Appendix
A3} \textbf{Corpus companion entry:} Claim C-187 \textbf{Audit Claim
ID:}
\texttt{ft-openai-o1-bioweapon-risk-2024-06-ai-double-use-reasoning-deepseek-claim-01}

\subsubsection{Appendix B17. Chapter 06 can ground the dual-use argument
in contemporaneous mainstream
reporti\ldots{}}\label{appendix-b17-chapter-06-can-ground-the-dual-use-argument-in-contemporaneous-mainstream-reporti}

\textbf{Claim statement:} That Chapter 06 can ground the dual-use
argument in contemporaneous mainstream reporting, not only in
retrospective critique. \textbf{Evidence status:} Documented Fact
\textbf{Source support:}
\hyperref[appendix-a3-openai-acknowledges-new-models-increase-risk-of-misuse-to-create-bioweapons]{Appendix
A3} \textbf{Corpus companion entry:} Claim C-188 \textbf{Audit Claim
ID:}
\texttt{ft-openai-o1-bioweapon-risk-2024-06-ai-double-use-reasoning-deepseek-claim-02}

\subsubsection{Appendix B18. Danger rhetoric around reasoning models
accompanied commercialization rather
than\ldots{}}\label{appendix-b18-danger-rhetoric-around-reasoning-models-accompanied-commercialization-rather-than}

\textbf{Claim statement:} That danger rhetoric around reasoning models
accompanied commercialization rather than halting deployment altogether.
\textbf{Evidence status:} Documented Fact \textbf{Source support:}
\hyperref[appendix-a3-openai-acknowledges-new-models-increase-risk-of-misuse-to-create-bioweapons]{Appendix
A3} \textbf{Corpus companion entry:} Claim C-189 \textbf{Audit Claim
ID:}
\texttt{ft-openai-o1-bioweapon-risk-2024-06-ai-double-use-reasoning-deepseek-claim-03}

\subsubsection{Appendix B19. There is direct technical pushback against
simplistic claims that frontier
reason\ldots{}}\label{appendix-b19-there-is-direct-technical-pushback-against-simplistic-claims-that-frontier-reason}

\textbf{Claim statement:} That there is direct technical pushback
against simplistic claims that frontier reasoning was trivially
replicated. \textbf{Evidence status:} Documented Fact \textbf{Source
support:}
\hyperref[appendix-a8-it-s-not-that-simple-an-analysis-of-simple-test-time-scaling]{Appendix
A8} \textbf{Corpus companion entry:} Claim C-243 \textbf{Audit Claim
ID:}
\texttt{its-not-that-simple-test-time-scaling-2025-06-ai-double-use-reasoning-deepseek-claim-01}

\subsubsection{Appendix B20. The paper argued simple test-time scaling
often reflects scaling down through
len\ldots{}}\label{appendix-b20-the-paper-argued-simple-test-time-scaling-often-reflects-scaling-down-through-len}

\textbf{Claim statement:} That the paper argued simple test-time scaling
often reflects scaling down through length constraints rather than true
scaling up. \textbf{Evidence status:} Documented Fact \textbf{Source
support:}
\hyperref[appendix-a8-it-s-not-that-simple-an-analysis-of-simple-test-time-scaling]{Appendix
A8} \textbf{Corpus companion entry:} Claim C-244 \textbf{Audit Claim
ID:}
\texttt{its-not-that-simple-test-time-scaling-2025-06-ai-double-use-reasoning-deepseek-claim-02}

\subsubsection{Appendix B21. The paper argued o1-like and
DeepSeek-R1-like systems differ from simple
replicat\ldots{}}\label{appendix-b21-the-paper-argued-o1-like-and-deepseek-r1-like-systems-differ-from-simple-replicat}

\textbf{Claim statement:} That the paper argued o1-like and
DeepSeek-R1-like systems differ from simple replications because they
learn to scale up test-time compute through reinforcement learning.
\textbf{Evidence status:} Documented Fact \textbf{Source support:}
\hyperref[appendix-a8-it-s-not-that-simple-an-analysis-of-simple-test-time-scaling]{Appendix
A8} \textbf{Corpus companion entry:} Claim C-245 \textbf{Audit Claim
ID:}
\texttt{its-not-that-simple-test-time-scaling-2025-06-ai-double-use-reasoning-deepseek-claim-03}

\subsubsection{Appendix B22. As of July 3, 2026, OpenAI's official API
pricing imposed a large price premium
o\ldots{}}\label{appendix-b22-as-of-july-3-2026-openai-s-official-api-pricing-imposed-a-large-price-premium-o}

\textbf{Claim statement:} That as of July 3, 2026, OpenAI's official API
pricing imposed a large price premium on \texttt{GPT-5.5\ Pro} relative
to \texttt{GPT-5.5}. \textbf{Evidence status:} Documented Fact
\textbf{Source support:}
\hyperref[appendix-a10-pricing-openai-api]{Appendix A10} \textbf{Corpus
companion entry:} Claim C-292 \textbf{Audit Claim ID:}
\texttt{openai-api-pricing-2026-06-ai-double-use-reasoning-deepseek-claim-01}

\subsubsection{Appendix B23. OpenAI's current API record treats top-tier
reasoning-capable access as
materiall\ldots{}}\label{appendix-b23-openai-s-current-api-record-treats-top-tier-reasoning-capable-access-as-materiall}

\textbf{Claim statement:} That OpenAI's current API record treats
top-tier reasoning-capable access as materially higher-cost
infrastructure rather than as a negligible increment over ordinary
flagship use. \textbf{Evidence status:} Documented Fact \textbf{Source
support:} \hyperref[appendix-a10-pricing-openai-api]{Appendix A10}
\textbf{Corpus companion entry:} Claim C-293 \textbf{Audit Claim ID:}
\texttt{openai-api-pricing-2026-06-ai-double-use-reasoning-deepseek-claim-02}

\subsubsection{Appendix B24. Chapter 06 can document reasoning-premium
monetization through a primary
pricing\ldots{}}\label{appendix-b24-chapter-06-can-document-reasoning-premium-monetization-through-a-primary-pricing}

\textbf{Claim statement:} That Chapter 06 can document reasoning-premium
monetization through a primary pricing source rather than only through
media summaries. \textbf{Evidence status:} Documented Fact
\textbf{Source support:}
\hyperref[appendix-a10-pricing-openai-api]{Appendix A10} \textbf{Corpus
companion entry:} Claim C-294 \textbf{Audit Claim ID:}
\texttt{openai-api-pricing-2026-06-ai-double-use-reasoning-deepseek-claim-03}

\subsubsection{Appendix B25. OpenAI publicly framed reasoning as a new
scaling layer based on additional
train\ldots{}}\label{appendix-b25-openai-publicly-framed-reasoning-as-a-new-scaling-layer-based-on-additional-train}

\textbf{Claim statement:} That OpenAI publicly framed reasoning as a new
scaling layer based on additional train-time and test-time compute.
\textbf{Evidence status:} Documented Fact \textbf{Source support:}
\hyperref[appendix-a1-learning-to-reason-with-llms]{Appendix A1}
\textbf{Corpus companion entry:} Claim C-310 \textbf{Audit Claim ID:}
\texttt{openai-learning-to-reason-with-llms-2024-06-ai-double-use-reasoning-deepseek-claim-01}

\subsubsection{Appendix B26. The company explicitly marketed o1 as a
major reasoning advance over GPT-4o in
ma\ldots{}}\label{appendix-b26-the-company-explicitly-marketed-o1-as-a-major-reasoning-advance-over-gpt-4o-in-ma}

\textbf{Claim statement:} That the company explicitly marketed
\texttt{o1} as a major reasoning advance over GPT-4o in math, coding,
and science-heavy evaluations. \textbf{Evidence status:} Documented Fact
\textbf{Source support:}
\hyperref[appendix-a1-learning-to-reason-with-llms]{Appendix A1}
\textbf{Corpus companion entry:} Claim C-311 \textbf{Audit Claim ID:}
\texttt{openai-learning-to-reason-with-llms-2024-06-ai-double-use-reasoning-deepseek-claim-02}

\subsubsection{Appendix B27. Chapter 06 can document reasoning as a
commercial and strategic product
category,\ldots{}}\label{appendix-b27-chapter-06-can-document-reasoning-as-a-commercial-and-strategic-product-category}

\textbf{Claim statement:} That Chapter 06 can document reasoning as a
commercial and strategic product category, not just a later media label.
\textbf{Evidence status:} Documented Fact \textbf{Source support:}
\hyperref[appendix-a1-learning-to-reason-with-llms]{Appendix A1}
\textbf{Corpus companion entry:} Claim C-312 \textbf{Audit Claim ID:}
\texttt{openai-learning-to-reason-with-llms-2024-06-ai-double-use-reasoning-deepseek-claim-03}

\subsubsection{Appendix B28. OpenAI itself paired stronger reasoning
capability with formal safety,
preparedne\ldots{}}\label{appendix-b28-openai-itself-paired-stronger-reasoning-capability-with-formal-safety-preparedne}

\textbf{Claim statement:} That OpenAI itself paired stronger reasoning
capability with formal safety, preparedness, and misuse-risk framing.
\textbf{Evidence status:} Documented Fact \textbf{Source support:}
\hyperref[appendix-a2-openai-o1-system-card]{Appendix A2} \textbf{Corpus
companion entry:} Claim C-316 \textbf{Audit Claim ID:}
\texttt{openai-o1-system-card-2024-06-ai-double-use-reasoning-deepseek-claim-01}

\subsubsection{Appendix B29. The company documented a training-data mix
combining public web-scale material,
p\ldots{}}\label{appendix-b29-the-company-documented-a-training-data-mix-combining-public-web-scale-material-p}

\textbf{Claim statement:} That the company documented a training-data
mix combining public web-scale material, proprietary partnership data,
and internal datasets. \textbf{Evidence status:} Documented Fact
\textbf{Source support:}
\hyperref[appendix-a2-openai-o1-system-card]{Appendix A2} \textbf{Corpus
companion entry:} Claim C-317 \textbf{Audit Claim ID:}
\texttt{openai-o1-system-card-2024-06-ai-double-use-reasoning-deepseek-claim-02}

\subsubsection{Appendix B30. Chapter 06 can ground the coexistence of
capability escalation and
safety-governa\ldots{}}\label{appendix-b30-chapter-06-can-ground-the-coexistence-of-capability-escalation-and-safety-governa}

\textbf{Claim statement:} That Chapter 06 can ground the coexistence of
capability escalation and safety-governance packaging in a primary
source. \textbf{Evidence status:} Documented Fact \textbf{Source
support:} \hyperref[appendix-a2-openai-o1-system-card]{Appendix A2}
\textbf{Corpus companion entry:} Claim C-318 \textbf{Audit Claim ID:}
\texttt{openai-o1-system-card-2024-06-ai-double-use-reasoning-deepseek-claim-03}

\subsubsection{Appendix B31. By November 23, 2023 there was serious
mainstream reporting explicitly linking
th\ldots{}}\label{appendix-b31-by-november-23-2023-there-was-serious-mainstream-reporting-explicitly-linking-th}

\textbf{Claim statement:} That by November 23, 2023 there was serious
mainstream reporting explicitly linking the board crisis to a letter
about an AI breakthrough. \textbf{Evidence status:} Documented Fact
\textbf{Source support:}
\hyperref[appendix-a4-openai-researchers-warned-board-of-ai-breakthrough-ahead-of-ceo-ouster-sources-say]{Appendix
A4} \textbf{Corpus companion entry:} Claim C-362 \textbf{Audit Claim
ID:}
\texttt{reuters-qstar-board-warning-2023-13-openai-governance-qstar-talent-exodus-claim-01}

\subsubsection{Appendix B32. Chapter 06 can describe the Q* thread as a
real reporting trail in the public
rec\ldots{}}\label{appendix-b32-chapter-06-can-describe-the-q-thread-as-a-real-reporting-trail-in-the-public-rec}

\textbf{Claim statement:} That Chapter 06 can describe the Q* thread as
a real reporting trail in the public record rather than a purely fringe
rumor. \textbf{Evidence status:} Documented Fact \textbf{Source
support:}
\hyperref[appendix-a4-openai-researchers-warned-board-of-ai-breakthrough-ahead-of-ceo-ouster-sources-say]{Appendix
A4} \textbf{Corpus companion entry:} Claim C-363 \textbf{Audit Claim
ID:}
\texttt{reuters-qstar-board-warning-2023-13-openai-governance-qstar-talent-exodus-claim-02}

\subsubsection{Appendix B33. The dossier should treat Q*-related causal
claims as reported and contested rathe\ldots{}
\{\#appendix-b33-the-dossier-should-treat-q-related-causal-claims-as-reported-and-contested-rathe\}}\label{appendix-b33.-the-dossier-should-treat-q-related-causal-claims-as-reported-and-contested-rathe-appendix-b33-the-dossier-should-treat-q-related-causal-claims-as-reported-and-contested-rathe}

\textbf{Claim statement:} That the dossier should treat Q*-related
causal claims as reported and contested rather than as settled fact.
\textbf{Evidence status:} Documented Fact \textbf{Source support:}
\hyperref[appendix-a4-openai-researchers-warned-board-of-ai-breakthrough-ahead-of-ceo-ouster-sources-say]{Appendix
A4} \textbf{Corpus companion entry:} Claim C-364 \textbf{Audit Claim
ID:}
\texttt{reuters-qstar-board-warning-2023-13-openai-governance-qstar-talent-exodus-claim-03}

\subsubsection{Appendix B34. By mid-July 2024 Reuters had published
serious reporting about an OpenAI
reasonin\ldots{}}\label{appendix-b34-by-mid-july-2024-reuters-had-published-serious-reporting-about-an-openai-reasonin}

\textbf{Claim statement:} That by mid-July 2024 Reuters had published
serious reporting about an OpenAI reasoning project under the codename
\texttt{Strawberry}. \textbf{Evidence status:} Documented Fact
\textbf{Source support:}
\hyperref[appendix-a5-exclusive-openai-working-on-new-reasoning-technology-under-code-name-strawberry]{Appendix
A5} \textbf{Corpus companion entry:} Claim C-365 \textbf{Audit Claim
ID:}
\texttt{reuters-strawberry-reasoning-2024-13-openai-governance-qstar-talent-exodus-claim-01}

\subsubsection{Appendix B35. The dossier can document a pre-release
public reporting trail for
reasoning-model\ldots{}}\label{appendix-b35-the-dossier-can-document-a-pre-release-public-reporting-trail-for-reasoning-model}

\textbf{Claim statement:} That the dossier can document a pre-release
public reporting trail for reasoning-model development before
\texttt{o1} entered the market. \textbf{Evidence status:} Documented
Fact \textbf{Source support:}
\hyperref[appendix-a5-exclusive-openai-working-on-new-reasoning-technology-under-code-name-strawberry]{Appendix
A5} \textbf{Corpus companion entry:} Claim C-366 \textbf{Audit Claim
ID:}
\texttt{reuters-strawberry-reasoning-2024-13-openai-governance-qstar-talent-exodus-claim-02}

\subsubsection{Appendix B36. Chapter 06 can connect later
reasoning-model branding to an earlier, serious
repo\ldots{}}\label{appendix-b36-chapter-06-can-connect-later-reasoning-model-branding-to-an-earlier-serious-repo}

\textbf{Claim statement:} That Chapter 06 can connect later
reasoning-model branding to an earlier, serious reporting record without
overstating certainty about every internal detail. \textbf{Evidence
status:} Documented Fact \textbf{Source support:}
\hyperref[appendix-a5-exclusive-openai-working-on-new-reasoning-technology-under-code-name-strawberry]{Appendix
A5} \textbf{Corpus companion entry:} Claim C-367 \textbf{Audit Claim
ID:}
\texttt{reuters-strawberry-reasoning-2024-13-openai-governance-qstar-talent-exodus-claim-03}

\subsubsection{Appendix B37. OpenAI's reasoning launch quickly triggered
open replication efforts focused on
t\ldots{}}\label{appendix-b37-openai-s-reasoning-launch-quickly-triggered-open-replication-efforts-focused-on-t}

\textbf{Claim statement:} That OpenAI's reasoning launch quickly
triggered open replication efforts focused on test-time scaling.
\textbf{Evidence status:} Documented Fact \textbf{Source support:}
\hyperref[appendix-a7-s1-simple-test-time-scaling]{Appendix A7}
\textbf{Corpus companion entry:} Claim C-368 \textbf{Audit Claim ID:}
\texttt{s1-simple-test-time-scaling-2025-06-ai-double-use-reasoning-deepseek-claim-01}

\subsubsection{Appendix B38. The paper proposed budget forcing as a
simple way to control test-time compute
by\ldots{}}\label{appendix-b38-the-paper-proposed-budget-forcing-as-a-simple-way-to-control-test-time-compute-by}

\textbf{Claim statement:} That the paper proposed budget forcing as a
simple way to control test-time compute by truncating or extending model
thinking. \textbf{Evidence status:} Documented Fact \textbf{Source
support:} \hyperref[appendix-a7-s1-simple-test-time-scaling]{Appendix
A7} \textbf{Corpus companion entry:} Claim C-369 \textbf{Audit Claim
ID:}
\texttt{s1-simple-test-time-scaling-2025-06-ai-double-use-reasoning-deepseek-claim-02}

\subsubsection{Appendix B39. The paper reported benchmark gains
substantial enough to challenge simplistic
fro\ldots{}}\label{appendix-b39-the-paper-reported-benchmark-gains-substantial-enough-to-challenge-simplistic-fro}

\textbf{Claim statement:} That the paper reported benchmark gains
substantial enough to challenge simplistic frontier-exclusivity
narratives. \textbf{Evidence status:} Documented Fact \textbf{Source
support:} \hyperref[appendix-a7-s1-simple-test-time-scaling]{Appendix
A7} \textbf{Corpus companion entry:} Claim C-370 \textbf{Audit Claim
ID:}
\texttt{s1-simple-test-time-scaling-2025-06-ai-double-use-reasoning-deepseek-claim-03}

\subsubsection{Appendix B40. Test-time compute scaling is a documented
technical mechanism rather than only
a\ldots{}}\label{appendix-b40-test-time-compute-scaling-is-a-documented-technical-mechanism-rather-than-only-a}

\textbf{Claim statement:} That test-time compute scaling is a documented
technical mechanism rather than only a commercial branding story.
\textbf{Evidence status:} Documented Fact \textbf{Source support:}
\hyperref[appendix-a6-scaling-llm-test-time-compute-optimally-can-be-more-effective-than-scaling-model-parameters]{Appendix
A6} \textbf{Corpus companion entry:} Claim C-376 \textbf{Audit Claim
ID:}
\texttt{snell-test-time-compute-2024-06-ai-double-use-reasoning-deepseek-claim-01}

\subsubsection{Appendix B41. The paper reported a compute-optimal
strategy improving test-time compute
efficie\ldots{}}\label{appendix-b41-the-paper-reported-a-compute-optimal-strategy-improving-test-time-compute-efficie}

\textbf{Claim statement:} That the paper reported a compute-optimal
strategy improving test-time compute efficiency by more than 4x over a
best-of-N baseline. \textbf{Evidence status:} Documented Fact
\textbf{Source support:}
\hyperref[appendix-a6-scaling-llm-test-time-compute-optimally-can-be-more-effective-than-scaling-model-parameters]{Appendix
A6} \textbf{Corpus companion entry:} Claim C-377 \textbf{Audit Claim
ID:}
\texttt{snell-test-time-compute-2024-06-ai-double-use-reasoning-deepseek-claim-02}

\subsubsection{Appendix B42. Under a FLOPs-matched evaluation the paper
found cases where a smaller model
usin\ldots{}}\label{appendix-b42-under-a-flops-matched-evaluation-the-paper-found-cases-where-a-smaller-model-usin}

\textbf{Claim statement:} That under a FLOPs-matched evaluation the
paper found cases where a smaller model using test-time compute
outperformed a model 14x larger. \textbf{Evidence status:} Documented
Fact \textbf{Source support:}
\hyperref[appendix-a6-scaling-llm-test-time-compute-optimally-can-be-more-effective-than-scaling-model-parameters]{Appendix
A6} \textbf{Corpus companion entry:} Claim C-378 \textbf{Audit Claim
ID:}
\texttt{snell-test-time-compute-2024-06-ai-double-use-reasoning-deepseek-claim-03}

\subsubsection{Appendix B43. Whether simple budget-forcing style
replications capture the same scaling
mechani\ldots{}}\label{appendix-b43-whether-simple-budget-forcing-style-replications-capture-the-same-scaling-mechani}

\textbf{Claim statement:} That whether simple budget-forcing style
replications capture the same scaling mechanism as o1-like or R1-like
reasoning remains contested. \textbf{Evidence status:} Disputed Fact
\textbf{Source support:}
\hyperref[appendix-a7-s1-simple-test-time-scaling]{Appendix A7},
\hyperref[appendix-a8-it-s-not-that-simple-an-analysis-of-simple-test-time-scaling]{Appendix
A8} \textbf{Corpus companion entry:} Claim C-434 \textbf{Audit Claim
ID:} \texttt{vector06-simple-replication-contested-claim-01}

\subsubsection{Appendix B44. The public record does not settle whether
Q* or any specific capability
breakthro\ldots{}}\label{appendix-b44-the-public-record-does-not-settle-whether-q-or-any-specific-capability-breakthro}

\textbf{Claim statement:} That the public record does not settle whether
Q* or any specific capability breakthrough directly triggered the
November 2023 OpenAI board rupture. \textbf{Evidence status:} Disputed
Fact \textbf{Source support:}
\hyperref[appendix-a4-openai-researchers-warned-board-of-ai-breakthrough-ahead-of-ceo-ouster-sources-say]{Appendix
A4} \textbf{Corpus companion entry:} Claim C-465 \textbf{Audit Claim
ID:} \texttt{vector13-qstar-causality-contested-claim-01}

\subsubsection{Appendix B45. Frontier labs are turning reasoning into a
premium political-economic category
by\ldots{}}\label{appendix-b45-frontier-labs-are-turning-reasoning-into-a-premium-political-economic-category-by}

\textbf{Claim statement:} That frontier labs are turning reasoning into
a premium political-economic category by coupling genuine capability
gains to danger rhetoric, controlled release, and tiered strategic
positioning. \textbf{Evidence status:} Hypothesis / Interpretation
\textbf{Source support:}
\hyperref[appendix-a2-openai-o1-system-card]{Appendix A2},
\hyperref[appendix-a3-openai-acknowledges-new-models-increase-risk-of-misuse-to-create-bioweapons]{Appendix
A3} \textbf{Corpus companion entry:} Claim C-433 \textbf{Audit Claim
ID:} \texttt{vector06-reasoning-premium-control-claim-01}

\subsubsection{Appendix B46. Current official pricing and launch records
from OpenAI and Anthropic show
reason\ldots{}}\label{appendix-b46-current-official-pricing-and-launch-records-from-openai-and-anthropic-show-reason}

\textbf{Claim statement:} That current official pricing and launch
records from OpenAI and Anthropic show reasoning capability being
commercialized through tiered plans, priority access, and premium API
pricing rather than released as an equal-access utility.
\textbf{Evidence status:} Hypothesis / Interpretation \textbf{Source
support:}
\hyperref[appendix-a9-chatgpt-plans-free-go-plus-pro-business-and-enterprise]{Appendix
A9}, \hyperref[appendix-a10-pricing-openai-api]{Appendix A10},
\hyperref[appendix-a11-plans-pricing-claude-by-anthropic]{Appendix A11},
\hyperref[appendix-a12-introducing-claude-4]{Appendix A12}
\textbf{Corpus companion entry:} Claim C-435 \textbf{Audit Claim ID:}
\texttt{vector06-tiered-reasoning-access-synthesis-claim-01}

\subsubsection{Appendix B47. Chapter 06 is strongest when it argues that
reasoning tiers package a real
comput\ldots{}}\label{appendix-b47-chapter-06-is-strongest-when-it-argues-that-reasoning-tiers-package-a-real-comput}

\textbf{Claim statement:} That Chapter 06 is strongest when it argues
that reasoning tiers package a real compute-intensive technical
mechanism inside a commercial and political access regime rather than
treating the whole phenomenon as fake. \textbf{Evidence status:}
Hypothesis / Interpretation \textbf{Source support:}
\hyperref[appendix-a6-scaling-llm-test-time-compute-optimally-can-be-more-effective-than-scaling-model-parameters]{Appendix
A6}, \hyperref[appendix-a2-openai-o1-system-card]{Appendix A2},
\hyperref[appendix-a9-chatgpt-plans-free-go-plus-pro-business-and-enterprise]{Appendix
A9}, \hyperref[appendix-a10-pricing-openai-api]{Appendix A10}
\textbf{Corpus companion entry:} Claim C-436 \textbf{Audit Claim ID:}
\texttt{vector06-ttc-real-but-packaged-claim-01}

\subsubsection{Appendix B48. OpenAI's governance rupture, safety
rhetoric, leadership churn, and
reasoning-mod\ldots{}}\label{appendix-b48-openai-s-governance-rupture-safety-rhetoric-leadership-churn-and-reasoning-mod}

\textbf{Claim statement:} That OpenAI's governance rupture, safety
rhetoric, leadership churn, and reasoning-model rollout are best read as
parts of one reorganization of authority rather than as isolated
episodes. \textbf{Evidence status:} Hypothesis / Interpretation
\textbf{Source support:}
\hyperref[appendix-a5-exclusive-openai-working-on-new-reasoning-technology-under-code-name-strawberry]{Appendix
A5} \textbf{Corpus companion entry:} Claim C-464 \textbf{Audit Claim
ID:} \texttt{vector13-governance-productization-synthesis-claim-01}

\subsubsection{Appendix B49. Open reasoning papers or public model
releases eliminate the politics of
compute\ldots{}}\label{appendix-b49-open-reasoning-papers-or-public-model-releases-eliminate-the-politics-of-compute}

\textbf{Claim statement:} That open reasoning papers or public model
releases eliminate the politics of compute concentration, distribution
control, or institutional gatekeeping around frontier AI.
\textbf{Evidence status:} Speculative Narrative Risk \textbf{Source
support:} \hyperref[appendix-a13-deepseek-v3-technical-report]{Appendix
A13},
\hyperref[appendix-a14-deepseek-r1-incentivizing-reasoning-capability-in-llms-via-reinforcement-learning]{Appendix
A14} \textbf{Corpus companion entry:} Claim C-432 \textbf{Audit Claim
ID:}
\texttt{vector06-open-publication-solves-control-overclaim-claim-01}

\subsubsection{Appendix B50. One secret breakthrough, letter, or hidden
internal discovery fully explains
the\ldots{}}\label{appendix-b50-one-secret-breakthrough-letter-or-hidden-internal-discovery-fully-explains-the}

\textbf{Claim statement:} That one secret breakthrough, letter, or
hidden internal discovery fully explains the OpenAI governance crisis,
later product strategy, and subsequent talent exodus by itself.
\textbf{Evidence status:} Speculative Narrative Risk \textbf{Source
support:}
\hyperref[appendix-a4-openai-researchers-warned-board-of-ai-breakthrough-ahead-of-ceo-ouster-sources-say]{Appendix
A4},
\hyperref[appendix-a5-exclusive-openai-working-on-new-reasoning-technology-under-code-name-strawberry]{Appendix
A5} \textbf{Corpus companion entry:} Claim C-466 \textbf{Audit Claim
ID:} \texttt{vector13-secret-trigger-overclaim-claim-01}

\subsection{Appendix C. Timeline
Slice}\label{appendix-c.-timeline-slice}

\subsubsection{Appendix C1. 2023-11-23: Reuters publishes a Q*-linked
reporting trail tied to the OpenAI board crisis
\{\#appendix-c1-2023-11-23-reuters-publishes-a-q-linked-reporting-trail-tied-to-the-openai-board-crisis\}}\label{appendix-c1.-2023-11-23-reuters-publishes-a-q-linked-reporting-trail-tied-to-the-openai-board-crisis-appendix-c1-2023-11-23-reuters-publishes-a-q-linked-reporting-trail-tied-to-the-openai-board-crisis}

\textbf{Event summary:} Reuters publishes a Q*-linked reporting trail
tied to the OpenAI board crisis \textbf{Event date:} 2023-11-23
\textbf{Source support:}
\hyperref[appendix-a4-openai-researchers-warned-board-of-ai-breakthrough-ahead-of-ceo-ouster-sources-say]{Appendix
A4} \textbf{Corpus companion entry:} Event T-033 \textbf{Audit Event
ID:} \texttt{event-reuters-qstar-report-2023-11-23}

\subsubsection{Appendix C2. 2024-07-15: Reuters reports on OpenAI's
Strawberry reasoning project before o1
release}\label{appendix-c2-2024-07-15-reuters-reports-on-openai-s-strawberry-reasoning-project-before-o1-release}

\textbf{Event summary:} Reuters reports on OpenAI's Strawberry reasoning
project before o1 release \textbf{Event date:} 2024-07-15 \textbf{Source
support:}
\hyperref[appendix-a5-exclusive-openai-working-on-new-reasoning-technology-under-code-name-strawberry]{Appendix
A5} \textbf{Corpus companion entry:} Event T-050 \textbf{Audit Event
ID:} \texttt{event-reuters-strawberry-2024-07-15}

\subsubsection{Appendix C3. 2024-08-06: Researchers publish a
compute-optimal test-time scaling paper for
LLMs}\label{appendix-c3-2024-08-06-researchers-publish-a-compute-optimal-test-time-scaling-paper-for-llms}

\textbf{Event summary:} Researchers publish a compute-optimal test-time
scaling paper for LLMs \textbf{Event date:} 2024-08-06 \textbf{Source
support:}
\hyperref[appendix-a6-scaling-llm-test-time-compute-optimally-can-be-more-effective-than-scaling-model-parameters]{Appendix
A6} \textbf{Corpus companion entry:} Event T-053 \textbf{Audit Event
ID:} \texttt{event-snell-test-time-compute-2024-08-06}

\subsubsection{Appendix C4. 2024-09-12: OpenAI launches o1-preview as a
reasoning-focused model
line}\label{appendix-c4-2024-09-12-openai-launches-o1-preview-as-a-reasoning-focused-model-line}

\textbf{Event summary:} OpenAI launches o1-preview as a
reasoning-focused model line \textbf{Event date:} 2024-09-12
\textbf{Source support:}
\hyperref[appendix-a1-learning-to-reason-with-llms]{Appendix A1}
\textbf{Corpus companion entry:} Event T-055 \textbf{Audit Event ID:}
\texttt{event-openai-o1-preview-2024-09-12}

\subsubsection{Appendix C5. 2024-09-13: Financial Times reports OpenAI
acknowledging higher bioweapon misuse risk for
o1}\label{appendix-c5-2024-09-13-financial-times-reports-openai-acknowledging-higher-bioweapon-misuse-risk-for-o1}

\textbf{Event summary:} Financial Times reports OpenAI acknowledging
higher bioweapon misuse risk for o1 \textbf{Event date:} 2024-09-13
\textbf{Source support:}
\hyperref[appendix-a3-openai-acknowledges-new-models-increase-risk-of-misuse-to-create-bioweapons]{Appendix
A3} \textbf{Corpus companion entry:} Event T-056 \textbf{Audit Event
ID:} \texttt{event-ft-openai-bioweapon-risk-2024-09-13}

\subsubsection{Appendix C6. 2024-12-05: OpenAI updates the o1 system
card with preparedness and training-data
details}\label{appendix-c6-2024-12-05-openai-updates-the-o1-system-card-with-preparedness-and-training-data-details}

\textbf{Event summary:} OpenAI updates the o1 system card with
preparedness and training-data details \textbf{Event date:} 2024-12-05
\textbf{Source support:}
\hyperref[appendix-a2-openai-o1-system-card]{Appendix A2} \textbf{Corpus
companion entry:} Event T-060 \textbf{Audit Event ID:}
\texttt{event-openai-o1-system-card-2024-12-05}

\subsubsection{Appendix C7. 2024-12-27: DeepSeek presents the
DeepSeek-V3 technical report on
arXiv}\label{appendix-c7-2024-12-27-deepseek-presents-the-deepseek-v3-technical-report-on-arxiv}

\textbf{Event summary:} DeepSeek presents the DeepSeek-V3 technical
report on arXiv \textbf{Event date:} 2024-12-27 \textbf{Source support:}
\hyperref[appendix-a13-deepseek-v3-technical-report]{Appendix A13}
\textbf{Corpus companion entry:} Event T-063 \textbf{Audit Event ID:}
\texttt{event-deepseek-v3-2024-12-27}

\subsubsection{Appendix C8. 2025-01-22: DeepSeek publishes its R1
reasoning paper and open release
claims}\label{appendix-c8-2025-01-22-deepseek-publishes-its-r1-reasoning-paper-and-open-release-claims}

\textbf{Event summary:} DeepSeek publishes its R1 reasoning paper and
open release claims \textbf{Event date:} 2025-01-22 \textbf{Source
support:}
\hyperref[appendix-a14-deepseek-r1-incentivizing-reasoning-capability-in-llms-via-reinforcement-learning]{Appendix
A14} \textbf{Corpus companion entry:} Event T-065 \textbf{Audit Event
ID:} \texttt{event-deepseek-r1-2025-01-22}

\subsubsection{Appendix C9. 2025-01-31: Researchers release s1 as an
open attempt to replicate test-time
scaling}\label{appendix-c9-2025-01-31-researchers-release-s1-as-an-open-attempt-to-replicate-test-time-scaling}

\textbf{Event summary:} Researchers release s1 as an open attempt to
replicate test-time scaling \textbf{Event date:} 2025-01-31
\textbf{Source support:}
\hyperref[appendix-a7-s1-simple-test-time-scaling]{Appendix A7}
\textbf{Corpus companion entry:} Event T-067 \textbf{Audit Event ID:}
\texttt{event-s1-simple-test-time-scaling-2025-01-31}

\subsubsection{Appendix C10. 2025-05-22: Anthropic launches Claude 4
with extended-thinking commercial
availability}\label{appendix-c10-2025-05-22-anthropic-launches-claude-4-with-extended-thinking-commercial-availability}

\textbf{Event summary:} Anthropic launches Claude 4 with
extended-thinking commercial availability \textbf{Event date:}
2025-05-22 \textbf{Source support:}
\hyperref[appendix-a12-introducing-claude-4]{Appendix A12}
\textbf{Corpus companion entry:} Event T-076 \textbf{Audit Event ID:}
\texttt{event-anthropic-claude-4-launch-2025-05-22}

\subsubsection{Appendix C11. 2025-07-19: A follow-up paper challenges
simplistic readings of simple test-time
scaling}\label{appendix-c11-2025-07-19-a-follow-up-paper-challenges-simplistic-readings-of-simple-test-time-scaling}

\textbf{Event summary:} A follow-up paper challenges simplistic readings
of simple test-time scaling \textbf{Event date:} 2025-07-19
\textbf{Source support:}
\hyperref[appendix-a8-it-s-not-that-simple-an-analysis-of-simple-test-time-scaling]{Appendix
A8} \textbf{Corpus companion entry:} Event T-082 \textbf{Audit Event
ID:} \texttt{event-its-not-that-simple-test-time-scaling-2025-07-19}

\subsection{Appendix D. Relevant
Entities}\label{appendix-d.-relevant-entities}

\subsubsection{Appendix D1. Anthropic}\label{appendix-d1-anthropic}

\textbf{Entity type:} company \textbf{Role in this paper:} Actor
relevant to 04\_data\_to\_models\_copyright\_weights,
06\_ai\_double\_use\_reasoning\_deepseek,
07\_export\_controls\_compute\_defense\_access,
08\_power\_networks\_legitimacy\_capture and others. \textbf{Relevant
sources in this paper:}
\hyperref[appendix-a11-plans-pricing-claude-by-anthropic]{Appendix A11},
\hyperref[appendix-a12-introducing-claude-4]{Appendix A12}
\textbf{Corpus companion entry:} Entity E-008 \textbf{Audit Entity ID:}
\texttt{org-anthropic}

\subsubsection{Appendix D2. DeepSeek}\label{appendix-d2-deepseek}

\textbf{Entity type:} company \textbf{Role in this paper:} Actor
relevant to 06\_ai\_double\_use\_reasoning\_deepseek,
07\_export\_controls\_compute\_defense\_access. \textbf{Relevant sources
in this paper:}
\hyperref[appendix-a13-deepseek-v3-technical-report]{Appendix A13},
\hyperref[appendix-a14-deepseek-r1-incentivizing-reasoning-capability-in-llms-via-reinforcement-learning]{Appendix
A14} \textbf{Corpus companion entry:} Entity E-013 \textbf{Audit Entity
ID:} \texttt{org-deepseek}

\subsubsection{Appendix D3. OpenAI}\label{appendix-d3-openai}

\textbf{Entity type:} company \textbf{Role in this paper:} Actor
relevant to 04\_data\_to\_models\_copyright\_weights,
06\_ai\_double\_use\_reasoning\_deepseek,
07\_export\_controls\_compute\_defense\_access,
08\_power\_networks\_legitimacy\_capture and others. \textbf{Relevant
sources in this paper:}
\hyperref[appendix-a1-learning-to-reason-with-llms]{Appendix A1},
\hyperref[appendix-a2-openai-o1-system-card]{Appendix A2},
\hyperref[appendix-a3-openai-acknowledges-new-models-increase-risk-of-misuse-to-create-bioweapons]{Appendix
A3},
\hyperref[appendix-a4-openai-researchers-warned-board-of-ai-breakthrough-ahead-of-ceo-ouster-sources-say]{Appendix
A4},
\hyperref[appendix-a5-exclusive-openai-working-on-new-reasoning-technology-under-code-name-strawberry]{Appendix
A5},
\hyperref[appendix-a9-chatgpt-plans-free-go-plus-pro-business-and-enterprise]{Appendix
A9}, \hyperref[appendix-a10-pricing-openai-api]{Appendix A10}
\textbf{Corpus companion entry:} Entity E-027 \textbf{Audit Entity ID:}
\texttt{org-openai}

\subsubsection{Appendix D4. Test-time compute
scaling}\label{appendix-d4-test-time-compute-scaling}

\textbf{Entity type:} technical\_mechanism \textbf{Role in this paper:}
Actor relevant to 06\_ai\_double\_use\_reasoning\_deepseek.
\textbf{Relevant sources in this paper:}
\hyperref[appendix-a6-scaling-llm-test-time-compute-optimally-can-be-more-effective-than-scaling-model-parameters]{Appendix
A6}, \hyperref[appendix-a7-s1-simple-test-time-scaling]{Appendix A7},
\hyperref[appendix-a8-it-s-not-that-simple-an-analysis-of-simple-test-time-scaling]{Appendix
A8} \textbf{Corpus companion entry:} Entity E-004 \textbf{Audit Entity
ID:} \texttt{concept-test-time-compute-scaling}

\subsection{Appendix E. Evidence
Boundaries}\label{appendix-e.-evidence-boundaries}

\subsubsection{Appendix E1. Disputed
Matters}\label{appendix-e1-disputed-matters}

\begin{itemize}
\tightlist
\item
  \hyperref[appendix-b43-whether-simple-budget-forcing-style-replications-capture-the-same-scaling-mechani]{Appendix
  B43}: That whether simple budget-forcing style replications capture
  the same scaling mechanism as o1-like or R1-like reasoning remains
  contested.
\item
  \hyperref[appendix-b44-the-public-record-does-not-settle-whether-q-or-any-specific-capability-breakthro]{Appendix
  B44}: That the public record does not settle whether Q* or any
  specific capability breakthrough directly triggered the November 2023
  OpenAI board rupture.
\end{itemize}

\subsubsection{Appendix E2. Interpretive
Boundaries}\label{appendix-e2-interpretive-boundaries}

\begin{itemize}
\tightlist
\item
  \hyperref[appendix-b45-frontier-labs-are-turning-reasoning-into-a-premium-political-economic-category-by]{Appendix
  B45}: That frontier labs are turning reasoning into a premium
  political-economic category by coupling genuine capability gains to
  danger rhetoric, controlled release, and tiered strategic positioning.
\item
  \hyperref[appendix-b46-current-official-pricing-and-launch-records-from-openai-and-anthropic-show-reason]{Appendix
  B46}: That current official pricing and launch records from OpenAI and
  Anthropic show reasoning capability being commercialized through
  tiered plans, priority access, and premium API pricing rather than
  released as an equal-access utility.
\item
  \hyperref[appendix-b47-chapter-06-is-strongest-when-it-argues-that-reasoning-tiers-package-a-real-comput]{Appendix
  B47}: That Chapter 06 is strongest when it argues that reasoning tiers
  package a real compute-intensive technical mechanism inside a
  commercial and political access regime rather than treating the whole
  phenomenon as fake.
\item
  \hyperref[appendix-b48-openai-s-governance-rupture-safety-rhetoric-leadership-churn-and-reasoning-mod]{Appendix
  B48}: That OpenAI's governance rupture, safety rhetoric, leadership
  churn, and reasoning-model rollout are best read as parts of one
  reorganization of authority rather than as isolated episodes.
\end{itemize}

\subsubsection{Appendix E3. Speculative Narrative
Risks}\label{appendix-e3-speculative-narrative-risks}

\begin{itemize}
\tightlist
\item
  \hyperref[appendix-b49-open-reasoning-papers-or-public-model-releases-eliminate-the-politics-of-compute]{Appendix
  B49}: That open reasoning papers or public model releases eliminate
  the politics of compute concentration, distribution control, or
  institutional gatekeeping around frontier AI.
\item
  \hyperref[appendix-b50-one-secret-breakthrough-letter-or-hidden-internal-discovery-fully-explains-the]{Appendix
  B50}: That one secret breakthrough, letter, or hidden internal
  discovery fully explains the OpenAI governance crisis, later product
  strategy, and subsequent talent exodus by itself.
\end{itemize}

\subsubsection{Appendix E4. Sensitive or Review-Worthy Note
Flags}\label{appendix-e4-sensitive-or-review-worthy-note-flags}

\begin{itemize}
\tightlist
\item
  \hyperref[appendix-a12-introducing-claude-4]{Appendix A12}. Companion
  note: Note N-054. Risk flag: low. Note focus: That on May 22, 2025
  Anthropic launched reasoning-oriented hybrid frontier models together
  with immediate commercial distribution across subscription plans, API
  access, and major cloud platforms.
\item
  \hyperref[appendix-a11-plans-pricing-claude-by-anthropic]{Appendix
  A11}. Companion note: Note N-055. Risk flag: low. Note focus: That as
  of July 3, 2026, Anthropic's official current pricing structure tied
  higher usage, priority access, and stronger feature access to premium
  paid plans.
\item
  \hyperref[appendix-a9-chatgpt-plans-free-go-plus-pro-business-and-enterprise]{Appendix
  A9}. Companion note: Note N-056. Risk flag: low. Note focus: That as
  of capture on July 3, 2026, OpenAI's official ChatGPT plan structure
  reserved the highest-end reasoning access for upper paid tiers rather
  than offering it uniformly across all users.
\item
  \hyperref[appendix-a14-deepseek-r1-incentivizing-reasoning-capability-in-llms-via-reinforcement-learning]{Appendix
  A14}. Companion note: Note N-057. Risk flag: medium. Note focus: That
  by January 2025 a non-U.S. lab was publicly claiming frontier-adjacent
  reasoning performance in a primary technical paper.
\item
  \hyperref[appendix-a13-deepseek-v3-technical-report]{Appendix A13}.
  Companion note: Note N-058. Risk flag: medium. Note focus: That by
  late 2024 there was already a competitor publicly claiming
  frontier-adjacent performance through a technical report.
\item
  \hyperref[appendix-a3-openai-acknowledges-new-models-increase-risk-of-misuse-to-create-bioweapons]{Appendix
  A3}. Companion note: Note N-059. Risk flag: medium. Note focus: That
  serious \texttt{T2} reporting documented OpenAI publicly acknowledging
  higher misuse risk at the same moment it was releasing stronger
  reasoning models.
\item
  \hyperref[appendix-a8-it-s-not-that-simple-an-analysis-of-simple-test-time-scaling]{Appendix
  A8}. Companion note: Note N-060. Risk flag: low. Note focus: That
  there is direct technical pushback against simplistic claims that
  frontier reasoning was trivially replicated.
\item
  \hyperref[appendix-a10-pricing-openai-api]{Appendix A10}. Companion
  note: Note N-061. Risk flag: low. Note focus: That as of July 3, 2026,
  OpenAI's official API pricing imposed a large price premium on
  \texttt{GPT-5.5\ Pro} relative to \texttt{GPT-5.5}.
\item
  \hyperref[appendix-a1-learning-to-reason-with-llms]{Appendix A1}.
  Companion note: Note N-062. Risk flag: medium. Note focus: That OpenAI
  publicly framed reasoning as a new scaling layer based on additional
  train-time and test-time compute.
\item
  \hyperref[appendix-a2-openai-o1-system-card]{Appendix A2}. Companion
  note: Note N-063. Risk flag: medium. Note focus: That OpenAI itself
  paired stronger reasoning capability with formal safety, preparedness,
  and misuse-risk framing.
\item
  \hyperref[appendix-a7-s1-simple-test-time-scaling]{Appendix A7}.
  Companion note: Note N-064. Risk flag: medium. Note focus: That
  OpenAI's reasoning launch quickly triggered open replication efforts
  focused on test-time scaling.
\item
  \hyperref[appendix-a6-scaling-llm-test-time-compute-optimally-can-be-more-effective-than-scaling-model-parameters]{Appendix
  A6}. Companion note: Note N-065. Risk flag: low. Note focus: That
  test-time compute scaling is a documented technical mechanism rather
  than only a commercial branding story.
\item
  \hyperref[appendix-a4-openai-researchers-warned-board-of-ai-breakthrough-ahead-of-ceo-ouster-sources-say]{Appendix
  A4}. Companion note: Note N-141. Risk flag: medium. Note focus: That
  by November 23, 2023 there was serious mainstream reporting explicitly
  linking the board crisis to a letter about an AI breakthrough.
\item
  \hyperref[appendix-a5-exclusive-openai-working-on-new-reasoning-technology-under-code-name-strawberry]{Appendix
  A5}. Companion note: Note N-142. Risk flag: medium. Note focus: That
  by mid-July 2024 Reuters had published serious reporting about an
  OpenAI reasoning project under the codename \texttt{Strawberry}.
\end{itemize}

\end{document}
